\begin{table}[htbp]
    \centering
    \small
    \caption{Bảng đối chiếu các khái niệm Cấu trúc rời rạc trong pipeline hệ thống}
    \begin{tabular}{|p{3.2cm}|p{4.2cm}|p{6.8cm}|}
    \hline
    \textbf{Chủ đề đề cương} & \textbf{Khái niệm toán học} & \textbf{Thành phần cài đặt trong hệ thống} \\ \hline
    \textbf{Sets \& Functions} & Tập hữu hạn, tập con, tích Descartes, ánh xạ vector & Phân loại 4 nhóm xe $C$, trục thời gian $T$, không gian $X \subseteq \mathbb{R}^{61}$ / $Y \subseteq \mathbb{R}^{60}$, hàm $f_{\text{WTI}}$, hàm làm mượt $s$ và ánh xạ hồi quy $F$. \\ \hline
    \textbf{Propositional Logic} & Mệnh đề chân trị $\mathbf{T}/\mathbf{F}$, phép hội ($\land$), tuyển ($\lor$), tương đương logic ($\equiv$) & Quy tắc cảnh báo vượt ngưỡng $3/5$ phút và hàm ước lượng thời gian đến ngưỡng ($\text{ETA}$) với toán tử $\min$. \\ \hline
    \textbf{Recursion} & Dãy số rời rạc, hệ thức truy hồi bậc cao & Trích xuất 16 đặc trưng trễ (Lags) ngắn, trung và dài hạn từ hai chuỗi quan sát $\text{total}$ và $\text{WTI}$. \\ \hline
    \textbf{Trees} & Cây nhị phân có gốc, đỉnh trong, đỉnh lá, chiều sâu cây & Cấu trúc cây quyết định trong mô hình HGB với ràng buộc $|V_L| \le 15$ dẫn đến giới hạn tổng số đỉnh $|V| \le 29$. \\ \hline
    \end{tabular}
\end{table}